\documentclass[11pt]{article}
\usepackage[a4paper,margin=27mm]{geometry}
\usepackage{amsmath,amssymb,mathtools}
\usepackage{fontspec,unicode-math}
\setmainfont{Latin Modern Roman}
\setmathfont{Latin Modern Math}
\usepackage{xurl,hyperref}
\hypersetup{colorlinks=true,urlcolor=blue,linkcolor=blue}
\newcommand{\sref}[2]{\hyperlink{source:#1:#2}{[S#2]}}
\newcommand{\cataloguescope}[1]{\paragraph{Mathematical question.} #1}
\begin{document}
\section{Conjecture 3.4 — Strong Four Exponentials Conjecture}\label{statement}

\subsection{Definitions and mathematical statement}

\paragraph{Shared notation.}$\mathbb N=\{1,2,\ldots\}$, and $\mathbb Z,\mathbb Q,\mathbb R,\mathbb C$ denote the integers, rationals, reals and complex numbers. Any other conventions are specified locally.


Write $\overline{\mathbb Q}$ for the algebraic numbers in $\mathbb C$ and set $\mathcal L=\{\lambda\in\mathbb C:e^\lambda\in\overline{\mathbb Q}^{\times}\}$. Define $\widetilde{\mathcal L}=\operatorname{span}_{\overline{\mathbb Q}}(\{1\}\cup\mathcal L)$: finite algebraic linear combinations of $1$ and arbitrary complex logarithms of nonzero algebraic numbers.
\paragraph{Statement recorded in the supplied source.}
If $x_1,x_2$ are linearly independent over $\overline{\mathbb Q}$ and $y_1,y_2$ are also linearly independent over $\overline{\mathbb Q}$, must at least one $x_i y_j$ ($1\le i,j\le2$) lie outside $\widetilde{\mathcal L}$? \sref{statement}{2}
\subsection{Short English statement}
Can both pairs of directions, each linearly independent over the algebraic numbers, have all four products in the algebraic span of one and logarithms of algebraic numbers?
\subsection{Sources}
\begin{description}
\item[\hypertarget{source:statement:1}{S1}] ULAM.AI and the UnsolvedMath Contributors, UnsolvedMath: A Curated Collection of Open Mathematics Problems, catalogue entry. CC BY 4.0. \url{https://huggingface.co/datasets/ulamai/UnsolvedMath}.
\item[\hypertarget{source:statement:2}{S2}] Michel Waldschmidt, Open Diophantine Problems, arXiv:math/0312440v2 (24 January 2004), Conjecture 3.4 and preceding definition of the logarithmic vector space. \url{https://arxiv.org/abs/math/0312440}.
\end{description}
\label{end:statement}
\end{document}
